% Build from this directory: latexmk main.tex
\documentclass[11pt,a4paper]{article}

\usepackage[T1]{fontenc}
\usepackage{lmodern}
\usepackage[margin=25mm]{geometry}

\begin{document}

\begin{titlepage}
  \centering
  \vspace*{\fill}
  {\LARGE\bfseries 3D Reconstruction of the UOB Innovation Hub\\
    Using Structure from Motion\par}
  \vspace{3em}

  \begin{tabular}{@{}ll@{}}
    \textbf{Name} & \textbf{Matriculation Number} \\[0.4em]
    % Replace the placeholder and add one row for each additional group member.
    [Full name] & [Matriculation number] \\
  \end{tabular}

  \vspace*{\fill}
  {\small\textbf{Generative AI disclosure.}
    OpenAI Codex was used to create the LaTeX report template.\par}
\end{titlepage}

\section{Introduction}
\label{sec:introduction}

\section{Methodology}
\label{sec:methodology}

\section{Experiments}
\label{sec:experiments}
% How to conduct experiments, e.g. settings, expected results

\subsection{Data acquisition}

We plan to take $N=25$ photographs of the UOB Innovation Hub and annotate
$M=8$ landmarks across the images. We will use a smartphone's main camera
and select landmarks on different facades or at different depths. We will
capture at least one pair with clear parallax and all eight landmarks
visible in both images. For the remaining photographs, we aim to keep at
least six of those landmarks visible to support camera registration.

\section{Discussion}
\label{sec:discussion}
% Real results, analysis

\section{Conclusion}
\label{sec:conclusion}

\section*{References}
% Use the heading text as the reference for this unnumbered section.
\begingroup
  \makeatletter
  \def\@currentlabel{References}
  \label{sec:references}
\endgroup

\end{document}
